\begin{afterword}
\summaryhead

The post consists of a single quotation, attributed to the psychologist and philosopher Eugene Gendlin, under the title ``You Can Face Reality.'' It has no commentary of its own.

The quotation makes five claims in seven short lines. What is true is already the case. Admitting it does not make it worse, and not being open about it does not make it go away. Because it is true, it is what a person actually has to deal with, while what is untrue is not there to be lived. People can bear the truth, because they are already bearing it.

The same passage closes the author's later post ``Avoiding Your Belief's Real Weak Points'' (October 2007). There he tells readers who are questioning a cherished belief to ``rehearse only this,'' and he cites Gendlin's book \textsc{Focusing}. On LessWrong the passage is known as the ``Litany of Gendlin.''

\responsehead

This is not an essay. It is a text for repetition. The author supplies a title, which turns the quotation into a promise, and line breaks, which turn a paragraph of prose into something that looks like a creed. The next time it appears, in ``Avoiding Your Belief's Real Weak Points,'' the reader is told to ``rehearse'' it. On the site it is called a litany, a word for a form of prayer. The form fits the use.

Its claims hold in easy cases and fail in hard ones, for the reasons given in the notes: owning up can make things worse, and the knowing is itself the new cost.

The quotation has also been moved. In Gendlin's book it follows the sentence ``It is not harmful to the other person if you look stupid or imperfect,'' in a part of the book about people helping each other. There the truth in question is the reader's own feeling, shown to a listener (my reading of a context I could see only in search snippets). The post gives no book or page, so the reader cannot recover that setting. It keeps the authority of the named source and drops the limits the source put on it.

The author has done this before. ``Twelve Virtues of Rationality'' rests its second virtue on another uprooted line, P.~C. Hodgell's ``That which can be destroyed by the truth should be.'' Hodgell has explained that in her novel \textsc{Seeker's Mask} (1994) the line is spoken by a character in ``the academic equivalent of a berserker fit---a ruthless drive to lay bare the truth, regardless of the cost,'' who is about to force another character to face facts ``at the worst possible time.'' In both cases the source placed the line in a setting that limits it, and the author presents it as a principle without one.

The author's own later rule condemns this use. ``Cached Thoughts'' (October 2007) warns against adopting another person's conclusion ready-made: ``if someone \emph{else} has already thought an idea through, you can save on computing power by caching their \emph{conclusion}---right?'' The litany is a conclusion packaged for use without examination, and its purpose is to be recited at the moment when examining would hurt.

The litany also tells its reciter who they are: someone who faces reality. It follows that whoever disagrees with them is someone who has flinched. That is an inference about how the text works; the post offers nothing to block it.

In short: a psychologist's remark about honesty with a friend, stripped of its context, set as a creed, and recommended for recitation by an author who elsewhere warns against repeating conclusions one has not examined.
\end{afterword}
